%% =====================================================================
%%  CAPITULO 5 - O FRAMEWORK SDD AGENTICO
%% =====================================================================

\chapter{O Framework SDD Agêntico}
\label{cap:framework}

Este capítulo descreve o processo proposto: sua arquitetura, os agentes, as
instruções permanentes, as habilidades de domínio, o fluxo adaptativo, os portões
de verificação, a definição de concluído, o protocolo de contexto e as condições
de parada. O framework é a variável independente do estudo.

\section{Princípios de projeto}
\label{sec:fw-principios}

O framework foi projetado a partir de oito princípios operacionais, adotados como
regras de conduta do executor:

\begin{enumerate}
	\item Planejar antes de editar.
	\item Verificar antes de assumir.
	\item Implementar pouco por vez.
	\item Testar imediatamente.
	\item Registrar decisões importantes.
	\item Nunca esconder incerteza.
	\item Nunca fabricar evidência.
	\item Nunca declarar validação física sem validação física.
\end{enumerate}

Os princípios 6 a 8 são os mais consequentes para o resultado deste trabalho.
Eles convertem uma postura epistemológica em regra operacional: a ausência de
dado não é uma falha a ser disfarçada, mas um estado a ser declarado. O item 8 é
específico ao domínio de sistemas embarcados do material de origem e estabelece
que nenhum nível de evidência pode ser promovido por analogia --- teste em
simulador não promove um resultado a validado em bancada.

\section{Motivação: lacunas do material de origem}
\label{sec:fw-lacunas}

O ponto de partida foi uma avaliação crítica de um conjunto prévio de artefatos
de processo. Essa avaliação identificou oito oportunidades de melhoria, listadas
no Quadro~\ref{qua:lacunas} com o tratamento dado a cada uma no framework
proposto.

\begin{quadro}[htb]
	\caption{Lacunas identificadas no material de origem e tratamento no framework}
	\label{qua:lacunas}
	\centering
	\footnotesize
	\begin{tabularx}{\textwidth}{@{}c p{5.4cm} X@{}}
		\toprule
		\textbf{\#} & \textbf{Lacuna identificada} & \textbf{Tratamento no framework proposto} \\
		\midrule
		L1 & Fluxo linear demais, com cadeia completa mesmo para alterações pequenas & Roteamento por complexidade e risco, com quatro trilhas (Seção~\ref{sec:fw-fluxo}) \\
		\addlinespace
		L2 & Ausência de um portão independente de verificação & Agente dedicado a responder se há evidência suficiente para declarar conclusão \\
		\addlinespace
		L3 & Economia de contexto não era regra operacional permanente & Conjunto de instruções permanentes com dez regras de contexto \\
		\addlinespace
		L4 & Falta de separação explícita entre estado real e memória da conversa & Protocolo de contexto com ordem de precedência de fontes \\
		\addlinespace
		L5 & Rastreabilidade definida em conceito, não operacionalizada & Matriz de rastreabilidade com formato fixo e regra de evidência obrigatória \\
		\addlinespace
		L6 & Risco de duplicação entre documentos de estado e de transferência & Definição de papel único por documento; transferência só quando há continuidade \\
		\addlinespace
		L7 & Ausência de mecanismo explícito de parada & Seis condições de parada obrigatórias \\
		\addlinespace
		L8 & Definição de concluído dispersa em vários documentos & Lista única de verificação de encerramento com quatorze itens \\
		\bottomrule
	\end{tabularx}
	\legend{Fonte: elaborado pelo autor (2026) a partir da avaliação crítica do material de origem.}
\end{quadro}

\section{Arquitetura em camadas}
\label{sec:fw-camadas}

O framework organiza-se em cinco camadas com responsabilidades distintas, como
ilustra a Figura~\ref{fig:camadas}. A separação é funcional: cada camada responde
a uma pergunta diferente, e a violação da fronteira entre camadas é a origem mais
comum de degradação de processo.

%% ---------------------------------------------------------------------
%%  FIGURA - Camadas do framework proposto (DG-02)
%% ---------------------------------------------------------------------
\begin{figure}[htb]
	\caption{Camadas do framework SDD agêntico e responsabilidade de cada uma}
	\label{fig:camadas}
	\centering
	\begin{tikzpicture}[
			font=\footnotesize,
			camada/.style={
					draw,
					rounded corners=2pt,
					minimum width=9.4cm,
					minimum height=1.0cm,
					align=center,
					fill=black!4
				}
		]
		\node[camada] (c1) at (0,0) {\textbf{Princípios}\\ constituição e regras permanentes};
		\node[camada, below=0.26cm of c1] (c2) {\textbf{Instruções}\\ contexto, mudanças, artefatos e domínio};
		\node[camada, below=0.26cm of c2] (c3) {\textbf{Habilidades}\\ conhecimento especializado sob demanda};
		\node[camada, below=0.26cm of c3] (c4) {\textbf{Agentes}\\ quinze papéis especializados};
		\node[camada, below=0.26cm of c4] (c5) {\textbf{Artefatos de controle}\\ especificação, decisões, tarefas, evidências, estado};

		\draw[->, thick] (-5.6,0.0) -- (-5.6,-5.5)
		node[midway, left, align=center, font=\scriptsize] {estabilidade\\ decrescente};
		\draw[->, thick] (5.9,-5.5) -- (5.9,0.0)
		node[midway, right, align=center, font=\scriptsize] {persistência\\ crescente};
	\end{tikzpicture}
	\legend{Fonte: elaborado pelo autor (2026) a partir da estrutura do repositório.}
\end{figure}


\textbf{Camada de princípios.} Contém o que não muda com a tarefa: a constituição
e as regras permanentes. É a camada mais estável e a de maior precedência.

\textbf{Camada de instruções.} Contém regras operacionais de aplicação ampla:
economia de contexto, controle de mudanças, artefatos e engenharia de domínio.
Instruções são carregadas por padrão de arquivo afetado, e não por decisão
discricionária.

\textbf{Camada de habilidades.} Contém conhecimento de domínio acionado sob
demanda: produção de especificação, desenvolvimento embarcado, artigo científico,
monografia e documentação de projeto. Habilidades não são carregadas
preventivamente.

\textbf{Camada de agentes.} Contém quinze papéis especializados, cada um com
entrada esperada, artefatos que pode ler, artefatos que pode alterar e portão de
saída.

\textbf{Camada de artefatos de controle.} Contém o estado do trabalho:
especificação, decisões, tarefas, estado, transferência e matriz de
rastreabilidade. É a única camada que persiste entre sessões.

\section{Agentes}
\label{sec:fw-agentes}

Os quinze agentes distribuem-se em quatro grupos funcionais, conforme o
Quadro~\ref{qua:agentes}.

\begin{quadro}[htb]
	\caption{Agentes do framework, por grupo funcional}
	\label{qua:agentes}
	\centering
	\footnotesize
	\begin{tabularx}{\textwidth}{@{}c p{3.4cm} X@{}}
		\toprule
		\textbf{\#} & \textbf{Agente} & \textbf{Responsabilidade} \\
		\midrule
		\multicolumn{3}{@{}l}{\textit{Grupo A --- Roteamento e planejamento}} \\
		\midrule
		00 & \texttt{task-router} & Classifica a solicitação, estima complexidade e risco e encaminha ao menor fluxo suficiente \\
		\addlinespace
		01 & \texttt{product-planner} & Define escopo, objetivos, restrições, riscos e critérios de sucesso \\
		\addlinespace
		02 & \texttt{requirements-engineer} & Redige requisitos funcionais e não funcionais, regras de negócio, casos de uso e critérios de aceitação \\
		\addlinespace
		03 & \texttt{architect} & Define componentes, interfaces, dados, tecnologias e decisões arquiteturais \\
		\addlinespace
		04 & \texttt{task-planner} & Decompõe requisitos e arquitetura em tarefas pequenas, ordenadas e verificáveis \\
		\midrule
		\multicolumn{3}{@{}l}{\textit{Grupo B --- Execução}} \\
		\midrule
		05 & \texttt{implementer} & Executa a tarefa aprovada seguindo requisitos, arquitetura e padrões do projeto \\
		\midrule
		\multicolumn{3}{@{}l}{\textit{Grupo C --- Verificação}} \\
		\midrule
		06 & \texttt{code-reviewer} & Avalia corretude, qualidade, manutenibilidade, testes e aderência à arquitetura \\
		\addlinespace
		07 & \texttt{test-engineer} & Define a estratégia de teste e cria testes de unidade, integração, sistema e ponta a ponta \\
		\addlinespace
		08 & \texttt{security-reviewer} & Avalia risco em código, arquitetura, dependências, configuração e tratamento de dados \\
		\addlinespace
		13 & \texttt{verification-gate} & Responde se há evidência suficiente para declarar a tarefa concluída \\
		\midrule
		\multicolumn{3}{@{}l}{\textit{Grupo D --- Encerramento e continuidade}} \\
		\midrule
		09 & \texttt{documentation} & Mantém a documentação técnica e de uso sincronizada \\
		\addlinespace
		10 & \texttt{integration-agent} & Valida a integração entre módulos, serviços, hardware e interfaces \\
		\addlinespace
		11 & \texttt{release-agent} & Verifica qualidade, documentação, versionamento e critérios de aceite antes da liberação \\
		\addlinespace
		12 & \texttt{project-orchestrator} & Coordena os demais agentes e preserva a rastreabilidade do roteamento à transferência \\
		\addlinespace
		14 & \texttt{handoff-manager} & Produz transferência compacta e verificável entre agentes ou sessões \\
		\bottomrule
	\end{tabularx}
	\legend{Fonte: elaborado pelo autor (2026) a partir de \texttt{.github/agents/}.}
\end{quadro}

A separação entre o agente que implementa (05) e os que verificam (06, 07, 13) é
deliberada e corresponde ao princípio clássico de independência entre construção e
verificação. O agente \texttt{code-reviewer} e o \texttt{test-engineer} avaliam;
o \texttt{verification-gate} decide. Concentrar as três funções em um único papel
reproduziria o problema que a separação pretende resolver.

\section{Instruções permanentes}
\label{sec:fw-instrucoes}

Quatro conjuntos de instruções são aplicados por padrão de arquivo afetado. Essa
forma de aplicação --- automática e baseada em caminho --- remove a decisão de
carregar ou não carregar da esfera discricionária do executor, o que é relevante
quando o executor é um agente.

\begin{quadro}[htb]
	\caption{Instruções permanentes do framework}
	\label{qua:instrucoes}
	\centering
	\footnotesize
	\begin{tabularx}{\textwidth}{@{}p{4.2cm} p{2.6cm} X@{}}
		\toprule
		\textbf{Instrução} & \textbf{Aplicação} & \textbf{Conteúdo} \\
		\midrule
		Economia de contexto & Todos os arquivos & Dez regras de inspeção incremental; proíbe varredura proativa; obriga leitura de estado antes de alterar \\
		\addlinespace
		Engenharia embarcada & Código C, C++, Rust, Python e assembly & Restrições de memória, tempo real, concorrência e segurança específicas do domínio \\
		\addlinespace
		Artefatos SDD & Especificações, documentação e arquivos de estado & Convenções de identificador (\texttt{FR}, \texttt{NFR}, \texttt{CA}, \texttt{T}, \texttt{ADR}) e regra de evidência obrigatória \\
		\addlinespace
		Controle de mudanças & Todos os arquivos & Escopo, menor alteração suficiente, proibição de alterar testes para fazê-los passar e verificação posterior obrigatória \\
		\bottomrule
	\end{tabularx}
	\legend{Fonte: elaborado pelo autor (2026) a partir de \texttt{.github/instructions/}.}
\end{quadro}

A instrução de controle de mudanças contém uma proibição que merece destaque:
é vedado alterar testes com o objetivo de fazê-los passar. Essa é a regra que
impede a degradação mais comum de processos de verificação --- a adaptação
retroativa do critério ao resultado obtido.

\section{Habilidades de domínio}
\label{sec:fw-skills}

Cinco habilidades concentram conhecimento especializado, carregado apenas quando a
tarefa correspondente está ativa:

\begin{itemize}
	\item \textbf{Criação de especificação} --- estrutura de especificação otimizada para consumo por agente, com distinção entre requisito, restrição e recomendação.
	\item \textbf{Desenvolvimento embarcado} --- especificação e implementação de firmware, com níveis de evidência e fluxo adaptativo.
	\item \textbf{Artigo científico em LaTeX} --- produção com rigor metodológico e referências verificáveis.
	\item \textbf{Monografia de engenharia em LaTeX} --- produção com conformidade ABNT, diagramas e projeto compilável.
	\item \textbf{Documentação de projeto} --- criação e manutenção do arquivo de apresentação do repositório.
\end{itemize}

O carregamento sob demanda é o mecanismo que resolve a tensão identificada na
Seção~\ref{sec:contexto-custo}: conhecimento especializado existe sem ocupar
contexto quando não é necessário.

\section{Fluxo adaptativo}
\label{sec:fw-fluxo}

O fluxo não aplica a cadeia completa a toda solicitação. O roteador classifica a
complexidade e o risco e seleciona a trilha mínima suficiente, como ilustra a
Figura~\ref{fig:fluxo-agentico}. Essa decisão responde diretamente à lacuna L1.

%% ---------------------------------------------------------------------
%%  FIGURA - Fluxo agentico adaptativo por complexidade (DG-01)
%% ---------------------------------------------------------------------
\begin{figure}[htb]
	\caption{Fluxo agêntico adaptativo: classificação por complexidade e risco}
	\label{fig:fluxo-agentico}
	\centering
	\begin{tikzpicture}[
			font=\footnotesize,
			no/.style={draw, rounded corners=2pt, minimum width=2.6cm, minimum height=0.85cm, align=center},
			dec/.style={draw, diamond, aspect=2.2, inner sep=1pt, align=center, font=\scriptsize},
			>=Stealth,
			fl/.style={->, thick}
		]

		\node[no] (req) at (0,0) {Solicitação};
		\node[no, below=0.7cm of req] (router) {Roteador};
		\node[dec, below=0.85cm of router] (risco) {complexidade\\ e risco?};

		\node[no, below left=1.5cm and 3.5cm of risco] (resp) {Resposta\\ com evidência};
		\node[no, below left=1.5cm and 0.9cm of risco] (quick) {Correção\\ local};
		\node[no, below right=1.5cm and 0.9cm of risco] (medio) {Recurso\\ médio};
		\node[no, below right=1.5cm and 3.5cm of risco] (grande) {Recurso\\ amplo};

		\draw[fl] (req) -- (router);
		\draw[fl] (router) -- (risco);
		\draw[fl] (risco) -- node[sloped, above, font=\scriptsize, pos=0.55] {pergunta} (resp);
		\draw[fl] (risco) -- node[sloped, above, font=\scriptsize] {simples} (quick);
		\draw[fl] (risco) -- node[sloped, above, font=\scriptsize] {médio} (medio);
		\draw[fl] (risco) -- node[sloped, above, font=\scriptsize, pos=0.55] {amplo} (grande);

		\node[no, below=1.0cm of medio] (impl) {Implementação};
		\draw[fl] (quick) |- (impl);
		\draw[fl] (medio) -- (impl);
		\draw[fl] (grande) |- (impl);

		\node[no, below=0.7cm of impl] (teste) {Testes};
		\node[no, below=0.7cm of teste] (revisao) {Revisão};
		\node[no, below=0.7cm of revisao] (gate) {Portão de\\ verificação};

		\draw[fl] (impl) -- (teste);
		\draw[fl] (teste) -- (revisao);
		\draw[fl] (revisao) -- (gate);
	\end{tikzpicture}
	\legend{Fonte: elaborado pelo autor (2026) a partir de \texttt{docs/agentic/WORKFLOW.md}.
	A trilha de liberação, aplicada após o portão de verificação, foi omitida para preservar a legibilidade.}
\end{figure}


O Quadro~\ref{qua:trilhas} descreve as quatro trilhas e o critério de elevação
entre elas.

\begin{quadro}[htb]
	\caption{Trilhas do fluxo adaptativo e critério de elevação}
	\label{qua:trilhas}
	\centering
	\footnotesize
	\begin{tabularx}{\textwidth}{@{}p{2.2cm} X p{3.4cm}@{}}
		\toprule
		\textbf{Trilha} & \textbf{Sequência de agentes} & \textbf{Critério de elevação} \\
		\midrule
		Pergunta & Resposta fundamentada, sem alteração de artefato & Necessidade de alterar qualquer arquivo \\
		\addlinespace
		Correção local & Roteador, executor, portão de verificação & Alteração de comportamento observável \\
		\addlinespace
		Recurso médio & Roteador, requisitos, executor, teste, revisão, verificação & Necessidade de decisão arquitetural \\
		\addlinespace
		Recurso amplo & Planejamento, requisitos, arquitetura, tarefas, execução, teste, revisão, segurança, verificação, documentação, integração & --- \\
		\addlinespace
		Liberação & Verificação, segurança, documentação, integração, liberação & --- \\
		\bottomrule
	\end{tabularx}
	\legend{Fonte: elaborado pelo autor (2026).}
\end{quadro}

\section{Portões de verificação}
\label{sec:fw-gates}

Seis portões pontuam o fluxo, conforme a Figura~\ref{fig:pipeline-gates} e o
Quadro~\ref{qua:gates}.

%% ---------------------------------------------------------------------
%%  FIGURA - Pipeline SDD com portoes de verificacao (DG-03)
%% ---------------------------------------------------------------------
\begin{figure}[htb]
	\caption{Pipeline SDD com os seis portões de verificação}
	\label{fig:pipeline-gates}
	\centering
	\begin{tikzpicture}[
			font=\footnotesize,
			etapa/.style={draw, rounded corners=2pt, minimum width=1.95cm, minimum height=0.9cm, align=center, fill=black!4},
			portao/.style={draw, fill=black!12, rounded corners=6pt, minimum width=0.85cm, minimum height=0.5cm, font=\scriptsize, align=center},
			>=Stealth,
			fl/.style={->, thick}
		]

		\node[etapa] (e1) at (0,0) {Estado\\ do repositório};
		\node[portao, right=0.35cm of e1] (g0) {G0};
		\node[etapa, right=0.35cm of g0] (e2) {Intenção\\ e requisito};
		\node[portao, right=0.35cm of e2] (g1) {G1};
		\node[etapa, right=0.35cm of g1] (e3) {Projeto e\\ decisões};
		\node[portao, right=0.35cm of e3] (g2) {G2};
		\node[etapa, below=1.1cm of e3] (e4) {Implementação\\ mínima};
		\node[portao, left=0.35cm of e4] (g3) {G3};
		\node[etapa, left=0.35cm of g3] (e5) {Evidência:\\ build e testes};
		\node[portao, left=0.35cm of e5] (g4) {G4};
		\node[etapa, left=0.35cm of g4] (e6) {Encerramento\\ e estado};
		\node[portao, left=0.35cm of e6] (g5) {G5};

		\draw[fl] (e1) -- (g0);
		\draw[fl] (g0) -- (e2);
		\draw[fl] (e2) -- (g1);
		\draw[fl] (g1) -- (e3);
		\draw[fl] (e3) -- (g2);
		\draw[fl] (g2) -- (e4);
		\draw[fl] (e4) -- (g3);
		\draw[fl] (g3) -- (e5);
		\draw[fl] (e5) -- (g4);
		\draw[fl] (g4) -- (e6);
		\draw[fl] (e6) -- (g5);
	\end{tikzpicture}
	\legend{Fonte: elaborado pelo autor (2026) a partir de \texttt{docs/agentic/WORKFLOW.md}.
	O portão G2 é aplicado apenas quando a complexidade exige decisão arquitetural.}
\end{figure}


\begin{quadro}[htb]
	\caption{Portões de verificação do framework}
	\label{qua:gates}
	\centering
	\footnotesize
	\begin{tabularx}{\textwidth}{@{}c p{3.2cm} X@{}}
		\toprule
		\textbf{Portão} & \textbf{Nome} & \textbf{Condição de passagem} \\
		\midrule
		G0 & Estado & Estado real do repositório inspecionado antes de qualquer alteração \\
		\addlinespace
		G1 & Intenção & Requisito e critério de aceitação identificados \\
		\addlinespace
		G2 & Projeto & Decisão arquitetural registrada quando a complexidade exige \\
		\addlinespace
		G3 & Implementação & Alteração mínima, rastreável e restrita ao escopo \\
		\addlinespace
		G4 & Evidência & Compilação, testes e validações aplicáveis executados \\
		\addlinespace
		G5 & Encerramento & Documentação, rastreabilidade, estado e riscos atualizados \\
		\bottomrule
	\end{tabularx}
	\legend{Fonte: elaborado pelo autor (2026) a partir de \texttt{docs/agentic/WORKFLOW.md}.}
\end{quadro}

O portão G1 é o que garante a precedência do critério sobre a implementação: sem
requisito e critério identificados, o fluxo não avança para a execução. O portão
G4 é o que impede a conclusão sem evidência.

\section{Definição de concluído}
\label{sec:fw-dod}

A definição de concluído consolida, em uma única lista de quatorze itens, o que
estava disperso em vários documentos de origem (lacuna L8):

1.~requisito relacionado identificado; 2.~critérios de aceitação verificados;
3.~estado do repositório inspecionado antes da alteração; 4.~alteração no menor
escopo necessário; 5.~compilação do alvo executada; 6.~testes executados;
7.~análise estática executada quando aplicável; 8.~consumo de memória e tempo
verificados quando aplicável; 9.~documentação atualizada; 10.~rastreabilidade
atualizada; 11.~arquivo de estado atualizado se houve mudança significativa;
12.~arquivo de transferência atualizado apenas se houver continuidade;
13.~riscos residuais registrados; 14.~nenhuma alteração não relacionada
incorporada.

Os itens 5, 6, 7 e 8 são ``quando aplicável''. Essa qualificação é deliberada: em
um artefato sem compilação --- como o produto deste estudo de caso, que é
interpretado diretamente pelo navegador --- não existe etapa de compilação, e
exigi-la produziria um item permanentemente não cumprido. A condição de
aplicabilidade precisa ser declarada, e não presumida.

\section{Protocolo de contexto}
\label{sec:fw-contexto}

O protocolo de contexto estabelece a ordem de precedência entre fontes de
informação:

\begin{enumerate}
	\item estado real do repositório e do sistema de controle de versão;
	\item princípios permanentes;
	\item especificação do recurso;
	\item registro de decisões arquiteturais;
	\item tarefas;
	\item código atual;
	\item testes e evidências;
	\item arquivos de estado e de transferência;
	\item histórico da conversa.
\end{enumerate}

A regra central é: \emph{o histórico da conversa pode fornecer intenção, mas nunca
deve substituir a leitura do estado atual do repositório}. Essa regra endereça a
lacuna L4 e é a defesa direta contra a falha mais insidiosa em fluxos executados
por agentes --- tratar uma versão lembrada de um arquivo como se fosse a versão
presente. A posição do histórico no último nível da hierarquia é a expressão
formal dessa defesa.

\section{Condições de parada}
\label{sec:fw-parada}

Seis condições obrigam a interrupção do fluxo e o pedido de decisão, em vez de
prosseguimento sob suposição:

\begin{enumerate}
	\item informação essencial de hardware ausente;
	\item duas especificações em conflito;
	\item necessidade de mudança arquitetural não aprovada;
	\item risco de perda de dados ou de persistência;
	\item impacto de segurança não especificado;
	\item teste obrigatório que não pode ser executado quando o resultado é necessário para liberar a tarefa.
\end{enumerate}

A segunda condição é a que produz mais consequências neste trabalho, porque uma
situação dessa natureza foi de fato encontrada: o repositório descreve um sistema
embarcado e contém, simultaneamente, um produto de página web. Diante do conflito,
a conduta adotada foi registrar a divergência e declarar a escolha, em vez de
produzir documentação sobre um sistema que não existe no repositório
(Seção~\ref{sec:delimitacao}).

\section{Matriz de rastreabilidade}
\label{sec:fw-rastreabilidade}

A matriz operacionaliza a rastreabilidade (lacuna L5) com sete colunas fixas:
requisito, critério, tarefa, código, teste, evidência e estado. Três regras a
governam: não marcar como aprovado sem evidência; acompanhar toda pendência de
justificativa e risco residual; e propagar toda alteração de requisito para
tarefa, teste e documentação.

Figura~\ref{fig:rastreabilidade} representa a cadeia de rastreabilidade aplicada
ao estudo de caso.

%% ---------------------------------------------------------------------
%%  FIGURA - Cadeia de rastreabilidade requisito-evidencia (DG-09)
%% ---------------------------------------------------------------------
\begin{figure}[htb]
	\caption{Cadeia de rastreabilidade do requisito à evidência}
	\label{fig:rastreabilidade}
	\centering
	\begin{tikzpicture}[
			font=\footnotesize,
			caixa/.style={draw, rounded corners=2pt, minimum height=0.85cm, align=center, fill=black!4, inner sep=4pt},
			>=Stealth,
			fl/.style={->, thick}
		]

		\node[caixa] (mem) at (0,0) {Solicitação\\ do produto};
		\node[caixa, right=0.6cm of mem] (req) {Requisitos\\ \texttt{FR}/\texttt{NFR}};
		\node[caixa, right=0.6cm of req] (ca) {Critérios\\ \texttt{CA}};
		\node[caixa, below=0.9cm of ca] (adr) {Decisões\\ \texttt{ADR}};
		\node[caixa, left=0.6cm of adr] (tar) {Tarefas\\ \texttt{T}};
		\node[caixa, left=0.6cm of tar] (cod) {Código\\ artefato};
		\node[caixa, below=0.9cm of cod] (evd) {Evidência\\ execução e inspeção};
		\node[caixa, right=0.6cm of evd] (res) {Resultado\\ PASS / PENDENTE};

		\draw[fl] (mem) -- (req);
		\draw[fl] (req) -- (ca);
		\draw[fl] (ca) -- (adr);
		\draw[fl] (adr) -- (tar);
		\draw[fl] (tar) -- (cod);
		\draw[fl] (cod) -- (evd);
		\draw[fl] (evd) -- (res);
	\end{tikzpicture}
	\legend{Fonte: elaborado pelo autor (2026). A ausência de critério rompe a cadeia entre requisito e
	evidência: é exatamente o caso dos dez requisitos marcados como \textbf{SEM CA} no Apêndice~\ref{ap:rastreabilidade}.}
\end{figure}


\section{Síntese do capítulo}
\label{sec:fw-sintese}

O framework proposto é um protocolo de controle: cinco camadas, quinze agentes,
quatro conjuntos de instruções, cinco habilidades, seis portões, catorze itens de
conclusão e seis condições de parada. Seus dois mecanismos distintivos são a
\emph{precedência do critério} sobre a implementação e a \emph{proibição de
evidência fabricada}. O Capítulo~\ref{cap:especificacao} mostra o primeiro em
operação; o Capítulo~\ref{cap:resultados}, o segundo.
